\documentclass[11pt]{article}
\usepackage[a4paper,margin=21mm]{geometry}
\usepackage[no-math]{fontspec}
\setmainfont{Latin Modern Roman}
\newfontfamily\CJKfont{Noto Serif CJK SC}
\XeTeXlinebreaklocale "zh"
\XeTeXlinebreakskip=0pt plus 1pt
\usepackage{amsmath,amssymb,amsthm,amscd,graphicx,color}
\usepackage{xurl}
\usepackage[unicode,hidelinks]{hyperref}
\allowdisplaybreaks
\setlength\emergencystretch{3em}
\numberwithin{equation}{section}

\numberwithin{equation}{section}

\newtheorem{Theorem}{Theorem}[section]
\newtheorem*{Theorem*}{Theorem}
\newtheorem{Corollary}[Theorem]{Corollary}
\newtheorem{Lemma}[Theorem]{Lemma}
\newtheorem{Proposition}[Theorem]{Proposition}
{ \theoremstyle{definition}
	\newtheorem{Definition}[Theorem]{Definition}
	\newtheorem{Note}[Theorem]{Note}
	\newtheorem{Example}[Theorem]{Example}
	\newtheorem{Remark}[Theorem]{Remark} }
	
\newcommand{\CC}{\ensuremath{\mathbb{C}}}
\newcommand{\RR}{\ensuremath{\mathbb{R}}}
\newcommand{\ZZ}{\ensuremath{\mathbb{Z}}}
\newcommand{\QQ}{\ensuremath{\mathbb{Q}}}
\newcommand{\NN}{\ensuremath{\mathbb{N}}}
\newcommand{\DD}{\ensuremath{\mathbb{D}}}
\newcommand{\md}{\mathrm{d}}
\newcommand{\me}{\mathrm{e}}
\newcommand{\mo}{\mathcal{O}}
\newcommand{\mi}{\mathrm{i}}

%\usepackage{float}
%\usepackage{changepage}
%\usepackage{setspace}
%\usepackage{graphicx}
%\usepackage{calc}

\newcommand{\tabincell}[2]{\begin{tabular}{@{}#1@{}}#2\end{tabular}}

\makeatletter\newlabel{Sect. 3}{{3}{}{Original source locator}{}{}}\makeatother
\begin{document}
\begin{center}\Large Complete asymptotic expansion of 2F2 with an integer parameter tending to negative infinity and nonresonant denominator parameters\end{center}
\noindent\textbf{Stable ID:} 1046\quad\textbf{Author:} Peng-Cheng Hang; Min-Jie Luo\par
\noindent\textbf{Work:} Asymptotics of the Humbert Function $\Psi_1$ for Two Large Arguments\par
\noindent\textbf{Source:} \url{https://sigma-journal.com/2024/074/}\par
\noindent\textbf{Version:} Official SIGMA 20 (2024) article 074, DOI 10.3842/SIGMA.2024.074; journal PDF and native TeX acquired 2026-09-30, exact bytes pinned by SHA256\par
\noindent\textbf{Rights:} CC BY 4.0. Authors retain copyright. \url{https://creativecommons.org/licenses/by/4.0/}. Reuse and typesetting adaptation; original language and mathematical content preserved.\par
\noindent\textbf{Difficulty (prior selection preserved):} {\CJKfont H1: The proof controls a hypergeometric transformation uniformly across five summation regimes, including indices near and beyond the moving negative parameter. Ordinary fixed-index Stirling estimates do not control that transition, so the author combines shifted-Gamma products, logarithmic cutoffs and factorial-tail estimates to justify the complete expansion. The nonresonance conditions c not a nonpositive integer, d not an integer and b-d not a nonnegative integer are retained.}\par
\medskip\noindent\textbf{Editorial boundary note (not author text):} Complete original Theorem4.3 for integer n tending to positive infinity, with all nonresonance conditions on c,d,b-d and z, every expansion term and the arbitrary-order remainder. The complete five-regime proof retains the fixed-index part, logarithmic cutoff, middle factorial ratio, moving-parameter transition products and full beyond-n factorial tail. Full original notation, Pochhammer-ratio proof, auxiliary factorial-series estimate and its precise textbook reference, complete shifted confluent-hypergeometric uniform-bound argument, exact Fields/Luke transformation formulas and preceding author context are included. The preceding sectorial lambda result and later general/other resonant variants are not additional counts. Intro105 explicitly positions Theorem4.3 among the paper main results; this is an independent complete expansion outcome, not a count for the preliminary estimate used in a separate Humbert application. Standard Stirling, textbook series comparison, pre-existing hypergeometric transformation and displayed prior bounds remain exact author prerequisites. Literal source equations, all five cases and indices are preserved; no proof generation, mathematical review or formula repair.\par
\noindent\textbf{External original reference Sect. 3:} Original independent Humbert-function two-large-variable application, source207–405, PDF4–8. Referenced only in the opening motivation of Section4; it is not required core for the selected full negative-integer-parameter expansion.\par
\medskip\hrule\medskip\noindent\textbf{Author text begins below.}\par
\setcounter{section}{1}\setcounter{equation}{2}
% Original source lines 107--205
\textbf{Notation.}
In this paper, the number $C$ generically denotes a positive constant independent of the parameter $n$, the index of summation $\ell$ and the variable $z$. Moreover, the generalized hypergeometric function $_pF_q$ is defined by (see, for example, \cite[p.~404]{NIST-Handbook})
	\begin{equation}\label{pFq definition}
		{}_pF_q\bigg[\begin{matrix}
			a_1,\dots,a_p\\
			b_1,\dots,b_q
		\end{matrix};z\bigg]
		\equiv
		{}_pF_q[
		a_1,\dots,a_p;
		b_1,\dots,b_q;z]
		:=\sum_{n=0}^{\infty}\frac{(a_1)_n\cdots(a_p)_n}{(b_1)_n\cdots(b_q)_n}\frac{z^n}{n!},
	\end{equation}
	where $a_1,\dots,a_p\in\CC$ and $b_1,\dots,b_q\in\CC\setminus\ZZ_{\leqslant 0}$.

\section{Preliminary lemmas}\label{Sect. 2}
In this section, we deduce three lemmas which will be used in the sequel. The first is a sharp bound for the ratio of Pochhammer symbols.
\begin{Lemma}\label{Lemma: ratio of Pochhammer symbol}
	If $a\in\CC$ and $b\in\CC\setminus\ZZ_{\leqslant 0}$, then
	\[\biggl|\frac{(a)_n}{(b)_n}\biggr|\leqslant Cn^{\operatorname{Re}(a-b)},\qquad n\in\ZZ_{>0}.\]
\end{Lemma}
\begin{proof}
	The proof for $a\in\ZZ_{\leqslant 0}$ is trivial. If $a\in\CC\setminus\ZZ_{\leqslant 0}$, Stirling's formula implies that
		\[\lim_{n\to\infty}n^{b-a}\frac{\Gamma(a+n)}{\Gamma(b+n)}=1,\]
		which concludes that $n^{b-a}\frac{(a)_n}{(b)_n}$ is bounded uniformly for $n\in\ZZ_{>0}$.
\end{proof}

The second is a simple estimate of the function $\Phi_a(x)$, which is the ``horizontal'' generating function of Stirling numbers of real order (see \cite[Section~8]{Butzer_Kilbas_Trujillo-2003}). More properties of $\Phi_a(x)$, containing the asymptotic behavior for fixed $x\in\RR$ and $a\to\pm\infty$, are studied in \cite[Section~3.2]{VanGorder-2009}.
\begin{Lemma}\label{Phi_a(x) estimate}
	Define
	\[\Phi_a(x):=\sum_{k=1}^{\infty}\frac{k^a}{k!}x^k,\qquad a,x\in\RR.\]
	Then $\Phi_a(x)\sim x^a\me^x$, $x\to +\infty$. Thus
	\[\Phi_a(x)\leqslant Kx^a\me^x,\qquad x\geqslant 1,\]
	where $K>0$ is a constant independent of $x$.
\end{Lemma}
\begin{proof}
	Take $a_n=\frac{n^a}{n!}$ and $b_n=\frac{1}{n!}$ in \cite[p.~12, Problem~72]{Polya_Szego-vol.2}.
\end{proof}

The third is a global estimate for the confluent hypergeometric function $_1F_1$.

\begin{Lemma}\label{lem-1F1}
	Let $a,c\in\CC$. Choose $N\in\ZZ_{>0}$ such that $\operatorname{Re}(c+N+1)>0$ and define
	\begin{equation}\label{G_l definition}
			G_{\ell}(z):={}_1F_1\bigg[\begin{matrix}
				a+\ell\\
				c+\ell
			\end{matrix};z\bigg],\qquad \ell\in\ZZ_{\geqslant 0}.
	\end{equation}
	Then for $\ell\geqslant N+1$,
 \begin{equation}\label{G_l inequality}
			|G_\ell(z)|\leqslant C\me^{\gamma|z|},
	\end{equation}
	where $\gamma>0$ is a constant independent of $\ell$, $N$ and $z$.
\end{Lemma}

\begin{proof}
	Since $G_{\ell}(0)=1$, we can assume that $z\ne 0$. Recall the inequality \cite[equation~(2.3)]{Joshi_Arya-1982}
	\begin{equation}\label{1F1 inequality-1}
		\biggl|{}_1F_1\bigg[\begin{matrix}
			a_0\\
			b_0
		\end{matrix};z\bigg]\biggr|\leqslant \cos\frac{\theta}{2}\cdot
		{}_1F_1\biggl[\begin{matrix}
			|a_0|\\
			|b_0|
		\end{matrix};|z|\sec\frac{\theta}{2}\bigg]
	\end{equation}
	with $\theta=\arg(b_0)\in (-\pi,\pi)$ and the inequality \cite[p. 37, equation~(3.5)]{Carlson-1966}
	\begin{equation}\label{1F1 inequality-2}
		\me^{\frac{a_0}{b_0}z}<{}_1F_1\biggl[\begin{matrix}
			a_0\\
			b_0
		\end{matrix};z\biggr]<1-\frac{a_0}{b_0}+\frac{a_0}{b_0}\me^z,\qquad b_0>a_0>0,\quad z\ne 0.
	\end{equation}
	But when $a_0\geqslant b_0>0$ and $z>0$, we have $\frac{(a_0)_k}{(b_0)_k}\leqslant \big(\frac{a_0}{b_0}\big)^k$ since $\frac{a_0+j}{b_0+j}$ decreases with respect to~${j\geqslant 0}$. Thus
	\begin{equation}\label{1F1 inequality-3}
		{}_1F_1\biggl[\begin{matrix}
			a_0\\
			b_0
		\end{matrix};z\biggr]\leqslant\me^{\frac{a_0}{b_0}z},\qquad a_0\geqslant b_0>0,\quad z>0.
	\end{equation}
	Recall $\operatorname{Re}(c+N+1)>0$ and note that $\gamma(\ell):=\max\bigl\{1,\frac{|a+\ell|}{|c+\ell|}\bigr\}$ is bounded uniformly for $\ell\in\ZZ_{\geqslant 0}$. Therefore, a combination of the inequalities \eqref{1F1 inequality-1}--\eqref{1F1 inequality-3} claims that for $\ell\geqslant N+1$,
	\[|G_\ell(z)|\leqslant 2\gamma(\ell)\me^{\sqrt{2}\gamma(\ell)|z|}\leqslant C\me^{\gamma|z|},\]
	where
	\[\gamma:=\sup_{\ell\in\ZZ_{\geqslant 0}}\sqrt{2}\gamma(\ell)\geqslant \sqrt{2}.\]
	This completes the proof.
\end{proof}

\begin{Remark}
	Lemma \ref{lem-1F1} can be easily generalized to the following form. Let $a_1,\dots,a_p,b_1,\dots$, $b_p\in\CC$, and let $N\in\ZZ_{>0}$ such that $\operatorname{Re}(b_j+N+1)>0$, $1\leqslant j\leqslant p$. Then for $\ell\geqslant N+1$,
		\begin{equation}\label{pFp estimate}
			\biggl|{}_pF_p\biggl[\begin{matrix}
				a_1+\ell,\dots,a_p+\ell\\
				b_1+\ell,\dots,b_p+\ell
			\end{matrix};z\biggr]\biggr|\leqslant C\me^{\gamma|z|},
		\end{equation}
		where $\gamma>0$ is a constant independent of $\ell$, $N$ and $z$.
\end{Remark}

\setcounter{section}{3}
% Original source lines 406--576
\section{Asymptotics of the generalized hypergeometric function}\label{Sect. 4}
Our derivation in Section~\ref{Sect. 3} depends on a rough estimate of $_2F_2$ for large $-n$. In this section, Nagel's approach is also used to explicitly establish the asymptotic behavior of $_2F_2$ for large parameters. We also present the asymptotic behavior of $_pF_q$ for large $-n$.

\subsection[Asymptotics of \_2F\_2]{Asymptotics of $\boldsymbol{_2F_2}$}
Let us examine the complete asymptotic expansions of $_2F_2$ for large parameters.

\begin{Theorem}\label{Thm: 2F2 large parameter--lambda}
	Let $\delta\in (0,\pi)$, $c\notin\ZZ_{\leqslant 0}$ and $z\ne 0$. Then for any positive integer $N$,
	\begin{equation}\label{2F2 large parameter--lambda}
		{}_2F_2\biggl[\begin{matrix}
			a,b+\lambda\\
			c,d+\lambda
		\end{matrix};z\biggr]=\sum_{k=0}^{N-1}\frac{(a)_k (d-b )_k}{(c)_k (d+\lambda )_k}\frac{ (-z )^k}{k!}{}_1F_1\biggl[\begin{matrix}
			a+k\\
			c+k
		\end{matrix};z\biggr]+\mo \bigl(\lambda^{-N} \bigr),
	\end{equation}
	where $\lambda\to\infty$ in the sector $|{\arg}(\lambda+d)|\leqslant \pi-\delta$.
\end{Theorem}

\begin{proof}
	If denoting
	\[v(z):={}_1F_1\biggl[\begin{matrix}
		a\\
		c
	\end{matrix};z\biggr],\qquad d_k:=\frac{(a)_k}{(c)_k k!},\]
	we can obtain
	\[{}_2F_2\biggl[\begin{matrix}
		a,b+\lambda\\
		c,d+\lambda
	\end{matrix};z\biggr]=\sum_{k=0}^{\infty}\frac{ (\lambda+b )_k}{ (\lambda+d )_k}d_k z^k.\]
	Note that the $k$-th derivative of $v(z)$ is given by \cite[p.~405, equation~(16.3.1)]{NIST-Handbook}
	\[v^{(k)}(z)=\frac{(a)_k}{(c)_k}
	{}_1F_1\biggl[\begin{matrix}
		a+k\\
		c+k
	\end{matrix};z\biggr].\]
	Applying Fields' result \cite[Theorem 3]{Fields-1967} yields
	\begin{gather}\label{2F2 transformation}
		{}_2F_2\biggl[\begin{matrix}
			a,b+\lambda\\
			c,d+\lambda
		\end{matrix};z\biggr]=\sum_{k=0}^{\infty}\frac{(d-b)_k}{(d+\lambda)_k}\frac{(-z)^k}{k!}v^{(k)}(z)=\sum_{k=0}^{\infty}\frac{(a)_k(d-b)_k}{(c)_k(d+\lambda)_k}\frac{(-z)^k}{k!}{}_1F_1\biggl[\begin{matrix}
			a+k\\
			c+k
		\end{matrix};z\biggr]
	\end{gather}
	with $\lambda+d\notin\ZZ_{\leqslant 0}$, and also gives the asymptotic expansion \eqref{2F2 large parameter--lambda}.
\end{proof}

\begin{Remark}
	~
	\begin{itemize}\itemsep=0pt
		\item[(1)] The series (\ref{2F2 transformation}) can be reformulated as follows:
		\begin{equation}\label{2F2 transformation reformulated}
			{}_2F_2\biggl[\begin{matrix}
				a,b\\
				c,d
			\end{matrix};z\biggr]=\sum_{k=0}^{\infty}\frac{(a)_k(d-b)_k}{(c)_k(d)_k}\frac{(-z)^k}{k!}{}_1F_1\biggl[\begin{matrix}
				a+k\\
				c+k
			\end{matrix};z\biggr],
		\end{equation}
		which is a specialization of Luke's result \cite[Section~9.1, equation~(34)]{Luke-1969}
		\begin{align}
					&{}_{p+1}F_{q+1}\biggl[\begin{matrix}
						a_1,\dots,a_p,b\\
						c_1,\dots,c_q,d
					\end{matrix};z\biggr]{}\nonumber\\
 & \qquad{}= \sum_{k=0}^{\infty}\frac{ (a_1 )_k\cdots (a_p )_k}{ (c_1 )_k\cdots (c_q )_k}\frac{(d-b)_k}{(d)_k}
 \frac{(-z)^k}{k!}{}_{p}F_{q}\biggl[\begin{matrix}
						a_1+k,\dots,a_p+k\\
						c_1+k,\dots,c_q+k
					\end{matrix};z\biggr],\label{pFp transformation}
\end{align}
			where $p\leqslant q$ with $z\in\CC$, or $p=q+1$ with $\operatorname{Re}(z)<\frac{1}{2}$.
		
		\item[(2)] Jur\v{s}\.enas \cite[Section~4.1, p.~69]{Jursenas-2014} mentioned an asymptotic formula
		\begin{equation}\label{Jursenas expansion-1}
			{}_2F_2\biggl[\begin{matrix}
				\alpha+1,\beta+\nu\\
				\delta,\beta+1+\nu
			\end{matrix};-\frac{1}{z}\biggr]\sim {}_1F_1\biggl[\begin{matrix}
				\alpha+1\\
				\delta
			\end{matrix};-\frac{1}{z}\biggr],\qquad \nu\in\ZZ_{\geqslant 0},\quad\nu\to \infty.
		\end{equation}
		without proof, whereas our Theorem \ref{Thm: 2F2 large parameter--lambda} gives a full asymptotic expansion of (\ref{Jursenas expansion-1}).
	\end{itemize}
\end{Remark}

The asymptotics of $_2F_2$ for $\lambda=-n$ is given below.
\begin{Theorem}\label{Thm: 2F2 large parameter--n}
	Let $c\notin\ZZ_{\leqslant 0}$, $d\notin\ZZ$, $b-d\notin\ZZ_{\geqslant 0}$ and $z\ne 0$. Then for any positive integer $N$,
	\begin{equation}\label{2F2 large parameter--n}
		{}_2F_2\biggl[\begin{matrix}
			a,b-n\\
			c,d-n
		\end{matrix};z\biggr]=\sum_{k=0}^{N-1}\frac{(a)_k(d-b)_k}{(c)_k(d-n)_k}\frac{(-z)^k}{k!}{}_1F_1\biggl[\begin{matrix}
			a+k\\
			c+k
		\end{matrix};z\biggr]+\mo\big(n^{-N}\big),
	\end{equation}
	where $n\to +\infty$ through integer values.
\end{Theorem}

\begin{proof}
	We shall follow Nagel's approach. For nonnegative integers $n$ and $\ell$, write
	\[F(z):={}_2F_2\biggl[\begin{matrix}
		a,b-n\\
		c,d-n
	\end{matrix};z\biggr],\qquad a_{\ell}(n):=\frac{(-1)^{\ell}}{(d-n)_{\ell}},\qquad g_{\ell}:=\frac{(a)_{\ell}(d-b)_{\ell}}{(c)_{\ell}\ell!}.\]
	Clearly, $\left|g_\ell\right|\leqslant C\ell^{\operatorname{Re}(\alpha)}$, where $\alpha:=a+d-b-c-1$. Moreover, \eqref{2F2 transformation reformulated} suggests that
	\begin{equation}\label{R(z) definition}
		R(z):=F(z)-G_0(z)=\sum_{\ell=1}^{\infty}a_\ell(n)g_\ell G_\ell(z)z^{\ell},
	\end{equation}
	where $G_\ell(z)$ is given by \eqref{G_l definition}.
	
	Take $n$ large, let $m=\lceil \log n\rceil$ and divide the series \eqref{R(z) definition} into five parts:
	\begin{equation}\label{five sums}
		R(z)=\sum_{1}^{N}+\sum_{N+1}^m+\sum_{m+1}^{n/2}+\sum_{n/2+1}^n+\sum_{n+1}^{\infty}=:S_0+S_1+S_2+S_3+S_4,
	\end{equation}
	where $N\in\ZZ_{>0}$ is chosen so that $\operatorname{Re}(c+N+1)>0$. Therefore, the inequality \eqref{G_l inequality} holds.
	
	Let us derive estimates for the sums in \eqref{five sums}.
	
	\textbf{Case 1.} $1\leqslant \ell\leqslant N$. Now $a_\ell(n)\sim n^{-\ell}$. Thus
	\[S_0=\sum_{\ell=1}^{N-1}a_\ell(n)g_\ell G_\ell(z)z^{\ell}+\mo\big(n^{-N}\big).\]
	
	\textbf{Case 2.} $N+1\leqslant\ell\leqslant m$. Now both $n$ and $n-\ell$ are large. The use of the identity $(z)_\ell=(-1)^{\ell}(1-z-\ell)_{\ell}$ gives
	\begin{equation}\label{a_l(n) identity}
		a_\ell(n)=\frac{1}{(1-d+n-\ell)_\ell}=\frac{\Gamma(1-d+n-\ell)}{\Gamma(1-d+n)}.
	\end{equation}
	By Stirling's formula, we get $a_\ell(n)\sim n^{-\ell}$ and thus
	\[|S_1|\leqslant C\mathrm{e}^{\gamma|z|}\sum_{\ell=N+1}^{m}\ell^{\operatorname{Re}(\alpha)}\biggl(\frac{|z|}{n}\biggr)^{\ell}.\]
	For $n$ large, $\ell^{\operatorname{Re}(\alpha)}\big(\frac{|z|}{n}\big)^{\ell}$ decreases with respect to $\ell\geqslant N+1$ and then
	\[|S_1|\leqslant C\mathrm{e}^{\gamma|z|}m(N+1)^{\operatorname{Re}(\alpha)}\biggl(\frac{|z|}{n}\biggr)^{N+1}=\mo\biggl(\frac{\log n}{n^{N+1}}\biggr).\]
	
	\textbf{Case 3.} $m+1\leqslant \ell\leqslant \frac{n}{2}$. Stirling's formula shows that
	\[a_\ell(n)=\frac{\Gamma(1-d+n-\ell)}{\Gamma(1-d+n)}\sim \biggl(1-\frac{\ell}{n}\biggr)^{-d}\frac{(n-\ell)!}{n!}.\]
	It follows from $\frac{1}{2}\leqslant 1-\frac{\ell}{n}<1$ that
	\[|a_\ell(n)|\leqslant C \frac{(n-\ell)!}{n!}=\frac{C}{(n-\ell+1)\cdots(n-1)n}\leqslant C\bigg(\frac{2}{n}\bigg)^{\ell}.\]
	As in Case 2, the monotonicity gives
	\[|S_2|\leqslant C\mathrm{e}^{\gamma|z|}\sum_{\ell=m+1}^{n/2}\ell^{\operatorname{Re}(\alpha)}\bigg(\frac{2|z|}{n}\bigg)^{\ell}\leqslant C\mathrm{e}^{\gamma|z|}nm^{\operatorname{Re}(\alpha)}\bigg(\frac{2|z|}{n}\bigg)^{m+1}=\mo\big(n^{-\frac{1}{2}\log n}\big).\]
	
	\textbf{Case 4.} $\frac{n}{2}+1\leqslant \ell\leqslant n$. Choose the least number $r\in\ZZ_{\geqslant 0}$ so that
	\[|(1-d+n-\ell)+r|\geqslant 4|z|>0.\]
	Note that $\frac{1}{|z+n|}$ is bounded uniformly for $n\geqslant 1$. Using \eqref{a_l(n) identity} gives
	\[|a_{\ell}(n)|=\prod_{j=0}^{\ell-1}\big|(1-d+n-\ell)+j\big|^{-1}\leqslant C(4|z|)^{r-\ell}.\]
	Thus
	\[|S_3|\leqslant C\mathrm{e}^{\gamma|z|}(4|z|)^r\sum_{\ell=n/2+1}^{n}\ell^{\operatorname{Re}(\alpha)}4^{-\ell}=\mo\big(n^{\operatorname{Re}(\alpha)+1}2^{-n}\big).\]
	
	\textbf{Case 5.} $\ell\geqslant n+1$. Now we can obtain $|a_\ell(n)|\leqslant C\frac{n^{\operatorname{Re}(d)}(\ell-n)^{1-\operatorname{Re}(d)}}{n!(\ell-n)!}$ from
	\[a_\ell(n)=\frac{(-1)^{\ell-n}}{(1-d)_n(d)_{\ell-n}}=\frac{(1)_n(1)_{\ell-n}}{(1-d)_n(d)_{\ell-n}}\frac{(-1)^{\ell-n}}{n!(\ell-n)!}.\]
	It follows that
	\begin{align*}
\begin{aligned}
		|S_4|&{} \leqslant C\mathrm{e}^{\gamma|z|}\frac{n^{\operatorname{Re}(d)}}{n!}\sum_{\ell=n+1}^{\infty}\ell^{\operatorname{Re}(\alpha)}(\ell-n)^{1-\operatorname{Re}(d)}\frac{|z|^{\ell}}{(\ell-n)!}\\
		&{} =C\mathrm{e}^{\gamma|z|}n^{\operatorname{Re}(d)}\frac{|z|^n}{n!}\sum_{j=1}^{\infty}(j+n)^{\operatorname{Re}(\alpha)}j^{1-\operatorname{Re}(d)}\frac{|z|^j}{j!}.
\end{aligned}
	\end{align*}
	Note that $\max\{n,j\}\leqslant j+n\leqslant 2\cdot\max\{n,j\}$. Then for $p,q\in\RR$, we have
	\begin{align*}
		\sum_{j=1}^{\infty}(j+n)^p j^q\frac{|z|^j}{j!}&{}\leqslant Cn^p\sum_{j=1}^{n}j^q\frac{|z|^j}{j!}+C\sum_{j=n+1}^{\infty}j^{p+q}\frac{|z|^j}{j!}\\
		&{} \leqslant Cn^{\max\{0,p\}}\sum_{j=1}^{\infty}j^{\max\{q,p+q\}}\frac{|z|^j}{j!},
	\end{align*}
	which shows that
	\[|S_4|\leqslant C\mathrm{e}^{\gamma|z|}n^{\max\{\operatorname{Re}(d),\operatorname{Re}(\alpha+d)\}}\frac{|z|^n}{n!}\sum_{j=1}^{\infty}j^{1-\operatorname{Re}(d)+\max\{0,\operatorname{Re}(\alpha)\}}\frac{|z|^j}{j!}.\]
	
	Now the asymptotic expansion \eqref{2F2 large parameter--n} follows from the estimates above.
\end{proof}
\begin{thebibliography}{99}
\bibitem[8]{Butzer_Kilbas_Trujillo-2003}
Butzer P.L., Kilbas A.A., Trujillo J.J., Stirling functions of the second kind
 in the setting of difference and fractional calculus, \href{https://doi.org/10.1081/NFA-120026366}{\textit{Numer. Funct.
 Anal. Optim.}} \textbf{24} (2003), 673--711.

\bibitem[9]{Carlson-1966}
Carlson B.C., Some inequalities for hypergeometric functions, \href{https://doi.org/10.2307/2035056}{\textit{Proc.
 Amer. Math. Soc.}} \textbf{17} (1966), 32--39.

\bibitem[12]{Fields-1967}
Fields J.L., Confluent expansions, \href{https://doi.org/10.2307/2004159}{\textit{Math. Comp.}} \textbf{21} (1967),
 189--197.

\bibitem[15]{Joshi_Arya-1982}
Joshi C.M., Arya J.P., Inequalities for certain hypergeometric functions,
 \href{https://doi.org/10.2307/2007475}{\textit{Math. Comp.}} \textbf{38} (1982), 201--205.

\bibitem[16]{Jursenas-2014}
Jur\v{s}\.enas R., On the definite integral of two confluent hypergeometric
 functions related to the {K}amp\'e de {F}\'eriet double series, \href{https://doi.org/10.1007/s10986-014-9227-y}{\textit{Lith.
 Math.~J.}} \textbf{54} (2014), 61--73, \href{https://arxiv.org/abs/1301.3039}{arXiv:1301.3039}.

\bibitem[19]{Luke-1969}
Luke Y.L., The special functions and their approximations. {V}ol.~{II},
 \textit{Math. Sci. Eng.}, Vol.~53, Academic Press, New York, 1969.

\bibitem[22]{NIST-Handbook}
Olver F.W.J., Olde~Daalhuis A.B., Lozier D.W., Schneider B.I., Boisvert R.F.,
 Clark C.W., Miller B.R., Saunders B.V., Cohl H.S., McClain M.A. (Editors),
 NIST digital library of mathematical functions, {R}elease~1.2.1 of
 2024-06-15, aviable at \url{https://dlmf.nist.gov/}.

\bibitem[24]{Polya_Szego-vol.2}
P\'olya G., Szeg\H{o} G., Problems and theorems in analysis. {V}ol.~{II}:
 {T}heory of functions, zeros, polynomials, determinants, number theory,
 geometry, \textit{Classics Math.}, \href{https://doi.org/10.1007/978-3-642-61905-2}{Springer}, Berlin, 1998.

\bibitem[25]{VanGorder-2009}
Van~Gorder R.A., Computation of certain infinite series of the form {$\sum
 f(n)n^k$} for arbitrary real-valued~{$k$}, \href{https://doi.org/10.1016/j.amc.2009.06.064}{\textit{Appl. Math. Comput.}}
 \textbf{215} (2009), 1209--1216.

\end{thebibliography}
\end{document}
